% STATUS: DRAFT
\chapter{Geometry from algebra}\label{ch:gelfand}
\skippable

\begin{physmotiv}
A classical phase space is a set of points on which observables are functions; the functions form a commutative algebra. Quantization replaces this algebra by a noncommutative one, and in that sense phase space itself becomes a ``noncommutative space''. Connes' noncommutative geometry takes this literally: a space \emph{is} its algebra of functions. This chapter explains why that is a theorem in the commutative case and what it suggests in general. Part IX is a self-contained detour; it connects back to Part IV through the modular flow (\cref{ch:thermal_time}).
\end{physmotiv}

\section{Gelfand--Naimark duality}

\begin{theorem}[Gelfand--Naimark]
Every commutative unital \cstar-algebra $\A$ is isomorphic to $C(X)$ for a compact Hausdorff space $X$, namely the space of characters (nonzero $^*$-homomorphisms $\A\to\C$, equivalently pure states) with the weak-$^*$ topology. Continuous maps $X\to Y$ correspond to $^*$-homomorphisms $C(Y)\to C(X)$ (contravariantly).
\end{theorem}

So compact spaces and commutative \cstar-algebras are the same thing, with points $=$ pure states (\cref{ch:states}: pure states of $C(X)$ are Dirac measures). Similarly, measure spaces correspond to commutative von Neumann algebras $\Linf(X,\mu)$, and the dictionary extends: \emph{noncommutative \cstar-algebras are noncommutative topological spaces; noncommutative von Neumann algebras are noncommutative measure spaces}.

\section{What is gained}
\begin{itemize}
\item \textbf{Singular quotients become algebras.} The quotient of a space by a group action that is not free (e.g.\ irrational rotation of the circle) is topologically pathological, but the crossed product $C(S^1)\rtimes\Z$ is a perfectly good \cstar-algebra, the ``noncommutative torus''.
\item \textbf{Phase space.} The Weyl algebra of \cref{ch:weyl} is the noncommutative torus algebra when $x,p$ are periodic; $\hbar$ is the noncommutativity parameter.
\item \textbf{Crossed products are spaces too.} The crossed product $\M\rtimes\R$ of \cref{ch:crossed} is the noncommutative version of the quotient of a ``space'' by an $\R$-action: for a measurable space with an action it is the groupoid algebra of the orbit equivalence relation. So the passage III$\to$II$_\infty$ is a geometric operation: ``forming the transversal''.
\end{itemize}

\begin{keyidea}
Commutative algebras are spaces. Noncommutative algebras are generalized spaces. Constructions on spaces (quotients, product, bundles) have algebraic counterparts (crossed products, tensor products, modules).
\end{keyidea}

\section*{Exercises}
\begin{exercise}
Show the characters of $C(X)$ are evaluations $f\mapsto f(x)$. (Hint: the kernel of a character is a maximal ideal; use Urysohn's lemma.)
\end{exercise}
\begin{exercise}
Show that $C(S^1)\rtimes_\theta\Z$ for the rotation by an angle $2\pi\theta$ is generated by unitaries $U,V$ with $UV=\ee^{2\pi\ii\theta}VU$ (noncommutative torus) and compare with the Weyl relations (\cref{ch:weyl}).
\end{exercise}
